\documentclass[aps,prl,twocolumn,showpacs,superscriptaddress]{revtex4-2}
\usepackage{amsmath,amssymb,booktabs,xcolor,hyperref,amsthm}
\newtheorem{theorem}{Theorem}

\begin{document}

\title{Carbon: The Topological Basis of Life\\
Why Carbon Alone Has $S_H = 1$: Maximum Hopf Bond Symmetry,
sp$^3$ Hybridisation, and the Bond Compression Theorem}

\author{Michel Robert Cabri\'e}
\affiliation{Independent Researcher, Victoria, Australia}
\email{ZeroFreeParameters@gmail.com}
\date{20 March 2026}

\begin{abstract}
We derive from BCT vacuum topology why carbon and carbon alone is 
the topological basis of life. Three results are established.
(1)~\textbf{Carbon Hopf Symmetry Theorem:} Carbon ($Z=6$, 
[He]$2s^22p^2$) has Hopf Bond Symmetry Index $S_H = 
\min(n_{\rm filled}, n_{\rm empty})/4 = 1.00$ --- the unique 
maximum among all $n=2$ elements. No other element has equal 
numbers of filled and open Hopf channels. $S_H=1$ means carbon 
can bond equally with itself, enabling arbitrarily long chains.
(2)~\textbf{sp$^3$ hybridisation} is the natural Hopf mode 
mixing at half-fill: equal admixture of $r_{\rm oct}$ (s) and 
$r_{\rm tet}$ (p) modes gives $\theta = \arccos(-1/3) = 
109.471^\circ$ (exact).
(3)~\textbf{Bond Compression Theorem:} each $\pi$ bond compresses 
C--C by $r_{\rm tet}$, giving C=C/C-C $= 1-r_{\rm tet}$ ($+2.0\%$) 
and C$\equiv$C/C-C $= 1-2r_{\rm tet}$ ($-0.5\%$). Life is 
carbon-based because BCT geometry has a unique half-fill symmetry 
point at the fourth element of the $n=2$ shell.
\end{abstract}

\maketitle

\section{Introduction}

Why is life based on carbon? Standard chemistry describes carbon's 
four bonds and self-chaining ability, but offers no explanation for 
why carbon specifically, among all light elements. BCT provides the 
answer: carbon occupies the unique half-fill point of the $n=2$ 
Hopf multiplet.

\section{The Hopf Bond Symmetry Index}

\begin{definition}
The Hopf Bond Symmetry Index: $S_H = \min(n_{\rm filled}, 
n_{\rm empty})/4$, where $n_{\rm filled}$ and $n_{\rm empty}$ are 
the filled and vacant outer-shell Hopf channels.
\end{definition}

\begin{center}
\begin{tabular}{lcccr}
\toprule
Element & $Z$ & Filled & Empty & $S_H$ \\
\midrule
Li & 3 & 1 & 7 & 0.25 \\
Be & 4 & 2 & 6 & 0.50 \\
B  & 5 & 3 & 5 & 0.75 \\
\textbf{C}  & \textbf{6} & \textbf{4} & \textbf{4} & \textbf{1.00} $\star$ \\
N  & 7 & 5 & 3 & 0.75 \\
O  & 8 & 6 & 2 & 0.50 \\
F  & 9 & 7 & 1 & 0.25 \\
Ne & 10 & 8 & 0 & 0.00 \\
\bottomrule
\end{tabular}
\end{center}

\begin{theorem}[Carbon Hopf Symmetry]
Carbon ($Z=6$) is the unique $n=2$ element with $S_H = 1.00$. 
It has 4 filled and 4 empty Hopf channels: equal capacity to 
donate and accept. $S_H = 1$ implies carbon can bond with itself: 
each C has open channels matching the filled channels of the 
other. This is why carbon chains are topologically self-consistent, 
why organic chemistry exists, and why life is carbon-based.
\end{theorem}

\section{Self-Bonding and Carbon Chains}

For a homoatomic chain A--A--A--$\ldots$, each interior atom must 
simultaneously donate and accept Hopf modes. This requires equal 
numbers of filled and empty channels: $S_H = 1$. Carbon satisfies 
this exactly. N ($S_H=0.75$) forms only short N--N bonds. O 
($S_H=0.50$) forms only O--O pairs. Si ($S_H=1$ at $n=3$) also 
chains, explaining silicon mineral structures.

\section{sp$^3$ Hybridisation at Half-Fill}

At $S_H=1$, carbon has equal occupancy of the $r_{\rm oct}$ mode 
(1 electron, s-character) and the $r_{\rm tet}$ modes (3 electrons, 
p-character). Half-fill forces equal mixing, producing four 
equivalent sp$^3$ hybrid channels at:
\begin{equation}
    \theta_{\rm sp^3} = \arccos(-1/3) = 109.471^\circ \quad\text{(exact)}
\end{equation}
sp$^3$ hybridisation is the topologically inevitable consequence 
of $S_H = 1$, not a choice made by carbon.

\section{The Bond Compression Theorem}

Each $\pi$ bond shares 2 additional Hopfions via $r_{\rm tet}$ 
modes perpendicular to the $\sigma$ axis, compressing the 
internuclear distance by $r_{\rm tet}$:

\begin{theorem}[Bond Compression]
For C--C bonds with $n_\pi$ additional $\pi$ bonds:
\begin{equation}
    \frac{L(n_\pi)}{L(0)} = 1 - n_\pi \times r_{\rm tet}
\end{equation}
\end{theorem}

\begin{center}
\begin{tabular}{llclr}
\toprule
Bond & $n_\pi$ & BCT ratio & Observed & Error \\
\midrule
C--C & 0 & $1$ (ref) & 1.540~\AA & --- \\
C=C  & 1 & $1-r_{\rm tet} = 0.8876$ & 1.340~\AA & $+2.0\%$ \\
C$\equiv$C & 2 & $1-2r_{\rm tet} = 0.7753$ & 1.200~\AA & $-0.5\%$ \\
\bottomrule
\end{tabular}
\end{center}

The triple bond ratio $1-2r_{\rm tet} = 0.7753$ matches observation 
to $-0.5\%$. The same $r_{\rm tet}=(\sqrt{6}-2)/4$ that gives the 
H$_2$O bond angle and the cosmic exhale also gives the C$\equiv$C 
compression.

\section{Conclusion}

Life is carbon-based because the $n=2$ Hopf multiplet has exactly 
one half-fill point, and that point is carbon. $S_H = 1.00$: equal 
filled and empty channels, self-bonding chains, exact tetrahedral 
geometry, bond compression by $r_{\rm tet}$ per $\pi$ bond. It 
could not be otherwise.

\begin{acknowledgments}
Derived in Melbourne, 20 March 2026, as part of the BCT chemistry 
series (Letters 78--82). The question ``why carbon?'' has a 
topological answer: $S_H = 1$.
\end{acknowledgments}

\begin{thebibliography}{99}
\bibitem{L78} M.~R.~Cabri\'e, \textit{BCT Letter 78}, Zenodo (2026).
\bibitem{L80} M.~R.~Cabri\'e, \textit{BCT Letter 80}, Zenodo (2026).
\bibitem{L81} M.~R.~Cabri\'e, \textit{BCT Letter 81}, Zenodo (2026).
\end{thebibliography}

\end{document}
